\section{System Architecture}
\label{sec:architecture}

Infon is built from four layers that compose without coupling: a cassette
substrate that stores infons as immutable, content-addressed files; a query DSL
that composes retrieval primitives over those files; a reasoner that converts
retrieved infons into calibrated Dempster--Shafer verdicts; and a schema
migration mechanism that evolves the ontology without reingesting documents.
Each design choice was validated by a reproducible probe; the key decisions
are described in the subsections below.

% -----------------------------------------------------------------------
\subsection{Cassette Substrate}
\label{sec:architecture:cassette}

\paragraph{File format.}
Each document is stored as a single \texttt{.inf} cassette.  The binary
layout is: an 8-byte magic number identifying the file type; a JSON header
containing metadata (schema reference, content hash, ingest timestamp); a
sequence of gzip-compressed record frames, one per infon; a Parquet index
footer; and a 16-byte trailer encoding the footer offset, allowing the index
to be retrieved with a single tail \texttt{GET} on S3.  Content identity is
defined by $\mathrm{sha256}(\text{text} \;\Vert\; \text{schema\_ref})$; two
ingests of the same document under the same schema produce identical cassette
identifiers, making delta ingest naturally idempotent.

\paragraph{Split indexes.}
Alongside each cassette, three Parquet shards are written:
\texttt{by\_triple}, \texttt{by\_time}, and \texttt{by\_anchor}.  These
indexes are never rewritten; appending new infons produces new shards that the
manifest accumulates.

\paragraph{Manifest pruner.}
The manifest maintains a cassette-level bounding box on anchor sets and time
ranges.  When a query arrives, the pruner eliminates cassettes whose bounding
box cannot intersect the query's anchor set or time window before any Parquet
file is opened.  Table~\ref{tab:pruner} summarises the measured result.

\begin{table}[ht]
  \centering
  \caption{Manifest pruner performance at 300 cassettes.}
  \label{tab:pruner}
  \begin{tabular}{ll}
    \toprule
    Mechanism & Result \\
    \midrule
    Parquet opens skipped (real workloads) & 7--16$\times$ fewer \\
    Parquet opens skipped (cache misses)   & 278$\times$ fewer \\
    Range gets per MCTS query (pruner + MCTS) & 1.4 average \\
    \bottomrule
  \end{tabular}
\end{table}

\paragraph{Time-travel via snapshot chain.}
Every ingest creates a new snapshot node in an append-only parent chain.
Calling \texttt{Manifest.load\_at(S)} returns a HEAD-as-of-snapshot-$S$ view
of the store, so historical queries are resolved without modifying any
cassette.  Post-migration time-travel is preserved: the snapshot chain holds
the pre-migration view at any prior snapshot identifier.

\paragraph{Delta ingest.}
\texttt{InfonStore.ingest()} is content-addressed by
$\mathrm{sha256}(\text{text} \;\Vert\; \text{schema\_ref})$.  If the cassette
already exists, ingest is a no-op; re-running the pipeline over the same
corpus costs nothing.

\paragraph{Schema migration.}
\texttt{SchemaFunctor(rename, merge, delete)} implements a Kan pushforward
(Section~\ref{sec:architecture:migration}).  Existing cassettes are rewritten
under the new ontology; old cassettes remain on disk and remain queryable at
their original snapshot identifiers.  Measured: migrating a 10-cassette,
600-infon store to a new schema takes 20\,ms, compared to re-extracting from
source documents, which takes on the order of seconds.  This is a 62$\times$
speedup.

% -----------------------------------------------------------------------
\subsection{Query DSL}
\label{sec:architecture:dsl}

The query DSL provides composable primitives over the cassette substrate.
Primitives are designed so that the manifest pruner and Parquet index pushdown
eliminate I/O before any infon is hydrated.

\paragraph{Triple-role filters.}
\texttt{where(s, p, o)} pins any combination of subject, predicate, and object;
unset arguments are treated as wildcards.  \texttt{mentioning(*a)} matches any
infon that contains anchor~$a$ in any role.

\paragraph{Polarity filters.}
\texttt{affirmed()} restricts to infons with polarity $+1$; \texttt{negated()}
restricts to polarity $-1$.  \texttt{contradicting()} constructs the
polarity-flipped counterpart of a retrieved infon, enabling direct refuter
lookup.

\paragraph{Temporal filters.}
\texttt{between(t\textsubscript{0}, t\textsubscript{1})}, \texttt{before(t)},
and \texttt{after(t)} filter on the ingest or publication timestamp stored in
the cassette header.  The manifest pruner intersects these windows against its
cassette-level time bounding boxes before opening any shard.

\paragraph{Schema-aware expansion.}
\texttt{expand\_hierarchy(schema)} rewrites an anchor query to include all
descendant entities in the ontology.  \texttt{run\_any([q\textsubscript{1},
\ldots, q\textsubscript{n}])} is logical OR across a list of queries, executed
as a single pass over the manifest.

\paragraph{Trajectory and aggregate primitives.}
\texttt{trajectory(anchor)} returns the time-ordered sequence of infons
involving an anchor, with NEXT edges derived at read time from the snapshot
chain.  \texttt{constraint(s, p, o)} computes corpus-level aggregates ---
\texttt{count\_by}, \texttt{first\_seen}, \texttt{last\_seen} --- via
aggregate pushdown: no infon records are hydrated, only the index shards are
read.

\paragraph{Pushdown efficiency.}
A representative compliance workload --- a 500-infon contradiction search ---
resolves to zero range gets: the query is answered entirely from index scans.
Time-travel snapshot resolution costs one additional file read regardless of
corpus size.

% -----------------------------------------------------------------------
\subsection{Reasoner}
\label{sec:architecture:reasoner}

The reasoner converts a set of retrieved infons into a calibrated
Dempster--Shafer mass $(m_S, m_R, m_U, m_\Theta)$ for a query triple.

\paragraph{Per-infon mass construction.}
Each infon contributes a base mass drawn from four sources:

\begin{enumerate}
  \item \textbf{Polarity.}  An affirmed infon ($\mathit{pol} = +1$) contributes
    mass to $m_S$; a negated infon ($\mathit{pol} = -1$) contributes mass to
    $m_R$.

  \item \textbf{Alignment.}  The SPLADE similarity score between the infon's
    anchor vector and the query anchor vector scales the committed mass; lower
    similarity shifts mass toward $m_\Theta$.

  \item \textbf{Distance.}  Infons reachable via longer hop-count chains
    contribute less committed mass; the remainder is added to $m_\Theta$,
    encoding greater ignorance for longer chains.

  \item \textbf{Confidence.}  The extractor's certainty score, derived from
    the SPLADE activation magnitude, further scales the committed mass.
\end{enumerate}

The canonical configuration used in all evaluations is \emph{top-1 cautious
fusion}: the single highest-scoring infon per claim is fused with $k = 2$
nearest neighbours and a cautious-weight parameter $c_w = 1.0$.  This
configuration was selected from an internal sweep over fusion parameters on
the synthetic held-out set.

\paragraph{Chain mass: conjunctive min/max (not Dempster's rule).}
When the reasoner traverses a multi-hop chain
$e_1 \to e_2 \to \cdots \to e_n$, it must aggregate the per-edge masses into
a single chain mass.  Dempster's combination rule is \emph{not} used for this
purpose.  The reason is fundamental: Dempster's rule amplifies $m_S$ as
additional edges are combined, because each combination renormalises away
conflict.  For a chain of assertions, the correct semantics is conjunctive ---
``all hops must hold'' --- which is naturally expressed by taking the
\emph{minimum} of supporting masses and the \emph{maximum} of refuting masses
along the chain:

\[
  m_S^{\text{chain}} = \min_{i} m_S^{(i)},
  \qquad
  m_R^{\text{chain}} = \max_{i} m_R^{(i)}.
\]

A negated edge anywhere in the chain propagates to the chain-level refutation
mass; a weak edge anywhere weakens the chain-level support.  Dempster's rule
would instead treat a chain of weak edges as strong combined evidence, which
is wrong for chained inference.

\paragraph{MCTS traversal.}
Monte Carlo Tree Search (MCTS) explores the infon hypergraph from a source
anchor to a target anchor.  Each expansion step follows a predicate-typed edge
to a neighbouring anchor; the chain mass is accumulated at each step.  The
sheaf GNN (Section~\ref{sec:background:sheaf_gnn}) serves as both the MCTS
prior --- its chain-verdict logit guides expansion --- and the terminal scorer
--- its final verdict head evaluates completed chains.

\paragraph{Connective-predicate auto-inference.}
MCTS expansion requires identifying which predicates are connective (i.e.,
link entities to entities) versus attributive (i.e., link entities to
property values).  Infon infers this automatically: a predicate is treated
as connective when its object-is-entity ratio exceeds 0.5 across the corpus.
Terminal anchors in chains additionally benefit from entity-widening, which
allows slight schema mismatches at chain endpoints.  No schema annotation is
required.

% -----------------------------------------------------------------------
\subsection{Schema Migration}
\label{sec:architecture:migration}

\paragraph{SchemaFunctor and Kan pushforward.}
Schema migration is expressed as a functor
$F \colon \mathcal{S}_{\text{old}} \to \mathcal{S}_{\text{new}}$ between
schema categories, with three primitive operations:
$\mathit{rename}$, $\mathit{merge}$, and $\mathit{delete}$.  The image of
each existing infon triple under $F$ is computed via the left Kan extension
of $F$ along the inclusion of the old schema, which extends $F$ to all
cassette records while preserving inter-entity relationships.

In practice, \texttt{SchemaFunctor(rename=\{...\}, merge=\{...\},
delete=\{...\})} rewrites every anchor in the store to its image under $F$.
Old cassettes are written to new files under the new schema; the original
files remain on disk.

\paragraph{Time-travel integrity.}
Because old cassettes are preserved and the snapshot chain records which
cassettes were visible at each snapshot identifier, queries issued against a
pre-migration snapshot still return the pre-migration view.  Migration does not
invalidate historical queries.

\paragraph{Performance.}
On a 10-cassette, 600-infon store, \texttt{store.migrate(functor,
``schema\_v2.json'')} completes in 20\,ms.  Re-extracting the same store from
source documents under the new schema would take on the order of seconds
(roughly 62$\times$ longer), because each document must be re-encoded by
SPLADE-tiny and re-parsed by the extraction pipeline.  The functor path
bypasses both steps.
